\section{2024-11-25 Verlust behaftet vs frei}

\startitemize[packed]
\item Welche gibt es?
\item Wo werden diese angewendet?
\item Eines vertieft erarbeiten
\stopitemize

\subsection{Welche Verfahren gibt es?}


\subsubsection{Bilder}

\startframedtext[width=\dimexpr\textwidth+\rightmarginwidth+\rightmargindistance\relax,frame=off,offset=none] % This spans textwidth + margin width
\startplacetable[location={here},title={}]
\setupTABLE[c][1][background=color,backgroundcolor=gray]
\setupTABLE[r][1][background=color,backgroundcolor=gray]
\bTABLE[option=stratch,offset=1mm]
\startCSV
[background=DiagonalRule] \DiagonalLabel{Verfahren}{Merkmal}, Verlust, Kompressionsrate, Anwendungsbereich, Dynamikbereich, Transparenz, Lizenz
png, frei, Besser bei einfachen Grafiken, "Buttons, Icons, Logos", bis zu 256 indizierte Farben, Ja, Frei
webp, beides, Bei hoher Komprimierung besser Bildqualität als JPEG, Web, 8 Bit Farbtiefe, Ja, BSD
gif (LZW), frei, abhängig von Farbpalette,"Computerspiele, Datenbanken, Banner, Kleine Animationene, In Chatts", 24-Bit,
\stopCSV
\eTABLE
\stopplacetable
\stopframedtext

\subsubsection{Videos}

\startplacetable[location={here},title={}]
\setupTABLE[c][1][background=color,backgroundcolor=gray]
\setupTABLE[r][1][background=color,backgroundcolor=gray]
\bTABLE[option=stratch,offset=1mm]
\startCSV
[background=DiagonalRule] \DiagonalLabel{Verfahren}{Merkmal}, Verlust, Kompressionsrate, Anwendungsbereich, Dynamikbereich, Lizenz
mp4, behaftet
MPEG-2, behaftet, mäßig
H.264, behaftet
WMV,
\stopCSV
\eTABLE
\stopplacetable

\subsubsection{Ton}

% \startframedtext[width=\dimexpr\textwidth+\rightmarginwidth+\rightmargindistance\relax,frame=off,offset=none] % This spans textwidth + margin width
\startplacetable[location={here},title={}]
\setupTABLE[c][1][background=color,backgroundcolor=gray]
\setupTABLE[r][1][background=color,backgroundcolor=gray]
\bTABLE[option=stratch,offset=1mm]
\startCSV
[background=DiagonalRule] \DiagonalLabel{Verfahren}{Merkmal}, Verlust, Ort, Übermittlung, Zugänglichkeit, Addressat, Intention
mp3, behaftet
flac, frei
wav, behaftet
AAC, behaftet
Apple Lossless, frei
\stopCSV
\eTABLE
\stopplacetable
% \stopframedtext

\subsubsection{Andere}

\startframedtext[width=\dimexpr\textwidth+\rightmarginwidth+\rightmargindistance\relax,frame=off,offset=none] % This spans textwidth + margin width
\startplacetable[location={here},title={}]
\setupTABLE[c][1][background=color,backgroundcolor=gray]
\setupTABLE[r][1][background=color,backgroundcolor=gray]
\bTABLE[option=stratch,offset=1mm]
\startCSV
[background=DiagonalRule] \DiagonalLabel{Verfahren}{Merkmal}, Verlust, Kompressionsrate, Anwendungsbereich, Geschwindigkeit, Lizens
zip, frei 
lz4, frei, , "Kompression: > 500 MB/s pro Kern, Decodieren: mehrer GB/s pro Kern", "Dateisysteme: ZFS und SquashFS, Linux Kernel", BSD 2-Clause und GPL-2.0-or-later
\stopCSV
\eTABLE
\stopplacetable
\stopframedtext


\subsection{LZ4}

\startitemize[packed]
\item Erscheinungsjahr: 2011
\item Ziel: Hohe Kompressions- und Dekompressionsgeschwindigkeit
\stopitemize

\subsubsection{Komprimiertes Datenformat}

Jeder komprimierter Block ist in Sequenzen aufgeteielt. Jede Sequenzen startet
mit einem {\bf Token}. So ein token ist ein Ein-Byte-Wert, welcher in zwei
4-Bits Felder aufgeteilt ist. Die ersten vier Bits geben die Länge der Literale

\subsubsection{Notes}





